\chapter{Message routing and dispatch internals}
\label{ch:dispatch}

\begin{aside}
Behind \code{std::visit} and \code{Cmd::map}, there is a real
dispatch machinery that turns variant types into jump tables and
threaded function pointers. This chapter opens the hood and walks
through what the compiler actually emits, how dispatch costs
compare across alternatives, and why \maya{}'s sum-type approach
beats both inheritance and \code{std::function} in practice.
\end{aside}

\section{The dispatch options C++ gives us}

Three mainstream ways to dispatch a message to a handler:

\begin{enumerate}[leftmargin=*]
  \item \textbf{Virtual functions.} Classic OOP. A vtable pointer
    per object, a vtable per class, one indirect call.
  \item \textbf{std::function or function pointer.} A type-erased
    callable; storage is heap-allocated for non-trivial
    closures.
  \item \textbf{Sum type + visit.} A discriminator integer plus
    inline storage; visit compiles to a switch.
\end{enumerate}

\maya{} uses option 3 throughout. The trade-offs are subtle but
real.

\subsection*{Virtual functions: the cost model}

\begin{lstlisting}
class Handler {
public:
    virtual void handle(Msg) = 0;
};

class Counter : public Handler {
public:
    void handle(Msg) override;
};

Handler* h = new Counter();
h->handle(msg);
\end{lstlisting}

The call \code{h->handle(msg)}:

\begin{itemize}[leftmargin=*]
  \item Load \code{h} (likely cached).
  \item Load vtable pointer at \code{*h} (one extra cache line).
  \item Load \code{handle} slot from vtable.
  \item Indirect call.
\end{itemize}

Costs: typically 2--5 cycles in the best case (vtable cached),
20+ cycles in the worst (vtable cold). Plus the heap allocation
for \code{new Counter()}.

\subsection*{std::function: the cost model}

\begin{lstlisting}
std::function<void(Msg)> handler = [](Msg m) { /* ... */ };
handler(msg);
\end{lstlisting}

Costs:

\begin{itemize}[leftmargin=*]
  \item If the captured state is small-object-optimised, two
    indirect calls (the \code{function} dispatches via vtable
    too).
  \item If the captured state exceeds SBO, one heap allocation
    plus the indirection.
  \item Inlining is rare; the compiler cannot prove
    \code{handler} points to a specific function.
\end{itemize}

Plus the heap allocation cost amortised across calls.

\subsection*{Sum type + visit: the cost model}

\begin{lstlisting}
std::variant<MsgA, MsgB, MsgC> msg = MsgA{};
std::visit([](auto& m) { /* ... */ }, msg);
\end{lstlisting}

Costs:

\begin{itemize}[leftmargin=*]
  \item Read the discriminator (one byte).
  \item Switch on the discriminator (jump table).
  \item Direct call to the matched branch.
  \item Inlining is common; the compiler sees every branch.
\end{itemize}

Costs: 1--3 cycles in the typical case. No heap allocation.

\section{What the compiler emits}

Let us look at the assembly.

\subsection*{Virtual call disassembly}

\begin{lstlisting}
; h->handle(msg)
mov rax, qword ptr [rdi]            ; load vtable ptr
mov rdx, qword ptr [rax]            ; load handle slot
jmp rdx                              ; indirect tail call
\end{lstlisting}

Three instructions, two memory loads, one indirect call. The
branch predictor must learn the target.

\subsection*{std::visit disassembly}

\begin{lstlisting}
; std::visit(f, v)
movzx eax, byte ptr [rdi + 16]      ; load discriminator
cmp eax, 3                           ; bounds check
jae default
jmp qword ptr [jump_table + rax*8]  ; jump-table dispatch

case_0:
    ; inlined body for MsgA
    ret
case_1:
    ; inlined body for MsgB
    ret
case_2:
    ; inlined body for MsgC
    ret
\end{lstlisting}

A jump table the size of the variant. Each case is inlined
where possible. The branch predictor learns the table entries.

\subsection*{The point}

The compiler can specialise the visit much more aggressively
than the virtual call. Inlining a variant case eliminates
function-call overhead; inlining a virtual call requires
devirtualisation, which is rare.

\section{Cmd::map: closure threading}

\code{Cmd::map} is more involved. It transforms a
\code{Cmd<ChildMsg>} into a \code{Cmd<ParentMsg>} via a
function. The implementation walks the variant and rebuilds it:

\begin{lstlisting}
template <class A, class B, class F>
Cmd<B> map(Cmd<A> cmd, F f) {
    return std::visit(overload{
        [&](None)      -> Cmd<B> { return Cmd<B>::none(); },
        [&](Quit)      -> Cmd<B> { return Cmd<B>::quit(); },
        [&](After<A> a) -> Cmd<B> {
            return Cmd<B>::after(a.delay, f(a.message));
        },
        [&](Task<A> t) -> Cmd<B> {
            auto fn = std::move(t.fn);
            auto f2 = std::move(f);
            return Cmd<B>::task([fn = std::move(fn), f = std::move(f2)]() mutable -> B {
                return f(fn());
            });
        },
        // ... one branch per alternative
    }, cmd.v);
}
\end{lstlisting}

\code{None}, \code{Quit}: no payload to transform; return the
same shape in the new type.

\code{After<A>}: transform the payload message via \code{f}.

\code{Task<A>}: build a new closure that runs the old task and
applies \code{f} to its result.

The cost: one allocation for \code{Task} (the closure storing
\code{f}); zero for the simpler alternatives.

\section{Why visit, not switch}

We could write \code{visit} by hand:

\begin{lstlisting}
switch (msg.index()) {
    case 0: handle(std::get<MsgA>(msg)); break;
    case 1: handle(std::get<MsgB>(msg)); break;
    case 2: handle(std::get<MsgC>(msg)); break;
}
\end{lstlisting}

This compiles to the same assembly as \code{std::visit} ---
but with two problems:

\begin{itemize}[leftmargin=*]
  \item Adding an alternative breaks the switch; the compiler
    does not warn unless you enable \code{-Wswitch}.
  \item \code{std::get<MsgA>} can throw if the variant is not
    \code{MsgA}; \code{visit}'s callable cannot.
\end{itemize}

\code{std::visit} is the type-safe spelling that the compiler
enforces. It is also the recipe the overload pattern targets.

\section{The exhaustiveness check}

When the variant grows, \code{std::visit} requires the visitor
to cover every alternative. If you forget one, the compiler
emits a clear error:

\begin{lstlisting}
error: no matching function for call to 'operator()' with 'MsgD'
\end{lstlisting}

For an overload-pattern visitor:

\begin{lstlisting}
error: 'overload<...>::operator()(MsgD)' is not callable
\end{lstlisting}

Either way: a new alternative without a corresponding handler
fails the build. You cannot ship a half-implemented message
variant.

\subsection*{When you want a catch-all}

Sometimes you want a default. An \code{auto} lambda matches any
alternative:

\begin{lstlisting}
std::visit(overload{
    [](MsgA a) { /* specific */ },
    [](MsgB b) { /* specific */ },
    [](auto)   { /* default */ }
}, msg);
\end{lstlisting}

The compiler picks the most-specialised overload for each
alternative; \code{auto} catches the rest. Use sparingly ---
it defeats exhaustiveness.

\section{Tagged union vs raw union}

\code{std::variant} is a tagged union: a discriminator plus
storage. A raw \code{union} has no discriminator; you must
track it yourself.

\begin{lstlisting}
union Msg {
    MsgA a;
    MsgB b;
    MsgC c;
};

// You must remember which one is active.
\end{lstlisting}

The trade-off:

\begin{itemize}[leftmargin=*]
  \item Variant: safer, slightly larger (one byte for tag plus
    alignment).
  \item Union: smaller, requires manual bookkeeping.
\end{itemize}

For one-off micro-optimisations a hand-rolled union can shave
bytes; for application code, the variant's safety wins every
time.

\section{Inlining and the optimiser}

Three factors determine how well the optimiser inlines visit:

\begin{enumerate}[leftmargin=*]
  \item \textbf{Visibility.} If the visitor and the visit call
    are in the same translation unit, the compiler can see all
    branches.
  \item \textbf{Constness.} A \code{const} variant is easier to
    reason about.
  \item \textbf{Branch count.} Above 10--12 branches, the
    compiler switches from inlined cases to a jump table; both
    are fast.
\end{enumerate}

In practice, \maya{}'s visit calls inline cleanly. The
generated code is competitive with hand-rolled switches.

\section{Pitfalls}

\begin{warning}
\textbf{Visit returning different types.} Each branch must
return a type the others convert to. If one returns
\code{int} and another \code{string}, the visit does not
compile. Wrap in \code{std::variant} or a base type.
\end{warning}

\begin{warning}
\textbf{Recursive visit.} Visiting from inside another visit
can be confusing. Capture the inner type clearly; write each
visit as its own named function and call them in sequence rather
than nesting visits at the call site.
\end{warning}

\begin{warning}
\textbf{Sharing state across visits.} The visitor lambda
captures by reference; the closure lives as long as the visit
returns. Do not store the visit-returned value beyond the
visitor's scope without copying.
\end{warning}

\begin{warning}
\textbf{Performance pessimism for small variants.}
\code{std::variant<int, double>} adds a byte. For tight inner
loops over millions of items, that byte can matter. Profile;
do not assume.
\end{warning}

\section{What the runtime emits per frame}

A typical \maya{} update loop has ten to twenty visit calls per
frame (one for the main message, one for each Cmd
interpretation, several for tree walks). The total dispatch
cost is microseconds, dominated by the actual work in each
branch, not by the dispatch itself.

If you ever profile a frame and the dispatch shows up high in
the flamegraph, the message variant is probably too large
(hundreds of alternatives) or the visit is in a tight loop
where it should not be. Move the visit outside the loop.

\section*{Workshop}

\begin{enumerate}[leftmargin=*]
  \item Write a small program with a 5-alternative variant and
    a \code{std::visit}. Compile with optimisation; inspect the
    assembly. Identify the jump table.
  \item Add a sixth alternative. Confirm the compiler errors
    on the visitor.
  \item Replace the visit with a \code{switch} on
    \code{msg.index()}. Confirm the assembly is identical.
\end{enumerate}

\begin{recap}
Sum types plus visit are the dispatch mechanism throughout
\maya{}. They compile to jump tables, inline cleanly, avoid
heap allocations, and enforce exhaustiveness. Virtual functions
and \code{std::function} have higher overhead and weaker static
guarantees. The choice is not stylistic; it is mechanical.
\end{recap}
